% ---------------------------------------------------------------------------
% Data generation: multislice 4D-STEM simulation + multislice electron
% ptychography. Style-matched to atomfind_methods.tex; these subsections
% precede sec:atomfind (they produce the object volume it operates on).
% Requires: amsmath. Citations \citep (natbib/biblatex); keys in
% sim_recon_refs.bib (+ Chen2021 from atomfind_refs.bib).
% ---------------------------------------------------------------------------
\subsection{Multislice 4D-STEM simulation}
\label{sec:sim}

\paragraph{Multislice forward model.}
Ground-truth four-dimensional scanning transmission electron microscopy (4D-STEM) data
are generated by the multislice method \citep{CowleyMoodie1957,Kirkland2010}, in which the
specimen electrostatic potential is partitioned into thin slices normal to the beam and the
electron wavefunction is propagated slice to slice by alternating transmission and Fresnel
propagation,
\begin{equation}
  \psi_{n+1} \;=\; \mathcal{P}_{\Delta z}\!\otimes\!\big[\,t_n(\mathbf{r})\,\psi_n(\mathbf{r})\,\big],
  \qquad t_n(\mathbf{r}) = \exp\!\big[\,i\,\sigma\, v_n(\mathbf{r})\,\big],
  \label{eq:multislice}
\end{equation}
with $t_n$ the transmission function of slice $n$, $v_n$ its projected potential, $\sigma$
the interaction constant, $\mathcal{P}_{\Delta z}$ the Fresnel propagator over slice
thickness $\Delta z$, and $\otimes$ convolution. We use the GPU implementation in abTEM
\citep{MadsenSusi2021}; projected potentials are built from the parameterisation of
\citet{LobatoVanDyck2014}, which is constrained to reproduce the correct elastic scattering
factors and electrostatic potentials, at $\Delta z = 2$\,\AA{} (reduced to $0.5$\,\AA{} for
the finely sampled reconstructions, well below the value at which the thin-slice
approximation of Eq.~\eqref{eq:multislice} biases the high-angle intensity). The specimen is
a PbTiO$_3$/SrTiO$_3$ ferroelectric labyrinth supercell ($70.0\times70.0\times47.9$\,\AA),
oriented with the beam along a $\langle 100\rangle$ axis and padded to a square $70$\,\AA{}
field of view, giving $\approx 70$\,\AA{} of material along the beam. The full box is
retained without cropping so that the laterally broadening exit wave does not alias at the
periodic boundaries.

\paragraph{Probe and scan.}
The convergent probe is defined by a $300$\,keV accelerating voltage, a $100$\,mrad
convergence semi-angle and a $20$\,\AA{} overfocus. It is raster-scanned over a $20$\,\AA{}
window at a step of $0.10$--$0.15$\,\AA, and the full diffraction intensity is recorded at
each position, i.e.\ a 4D-STEM dataset \citep{Ophus2019}. The large convergence angle and
the resulting strong overlap of neighbouring probe positions oversample the object and are
what render the phase problem well-conditioned for the ptychographic inversion of
Sec.~\ref{sec:recon} \citep{RodenburgFaulkner2004}.

\paragraph{Thermal diffuse scattering.}
Thermal scattering is included in the frozen-phonon approximation
\citep{LoaneXuSilcox1991}: the recorded intensity is the incoherent average over
$N_\mathrm{ph}=16$ statistically independent snapshots of the lattice,
\begin{equation}
  I(\mathbf{k}) \;=\; \frac{1}{N_\mathrm{ph}}\sum_{j=1}^{N_\mathrm{ph}}
  \big|\Psi_j(\mathbf{k})\big|^2,
  \label{eq:fp}
\end{equation}
each snapshot displacing every atom by an independent Gaussian draw. Amplitudes are set
\emph{per species} from tabulated room-temperature isotropic Debye--Waller factors $B$
\citep{Peng1996}: the one-dimensional root-mean-square displacement is
$u_{1\mathrm{D}}=\sqrt{B/8\pi^2}$, and because abTEM parameterises the three-dimensional
displacement magnitude we pass $\sigma=\sqrt{3B/8\pi^2}$. The heavy A-site lead and the
oxygen vibrate roughly twice as far as titanium ($B_\mathrm{Pb}=0.90$, $B_\mathrm{O}=0.80$,
$B_\mathrm{Sr}=0.55$, $B_\mathrm{Ti}=0.45$\,\AA$^2$), so a single averaged displacement would
misweight the diffuse background of precisely the light species we most wish to detect; the
per-species treatment is therefore not cosmetic for the oxygen problem of
Sec.~\ref{sec:atomfind}. The convergence of Eq.~\eqref{eq:fp} in $N_\mathrm{ph}$ was checked
against a doubled configuration count.

\paragraph{Detector and dose.}
The full $\approx 200$\,mrad diffraction plane is recorded and binned $4\times4$ to
$356\times356$ pixels, keeping the zero-order disc centred so that the reciprocal-space
calibration required by the reconstruction is exact. abTEM flux-normalises each pattern; to
impose a finite electron budget we rescale to an incident fluence and draw every pixel from
a Poisson distribution with mean equal to the noiseless intensity scaled to $D\,s^2$
electrons per pattern, where $D$ is the dose in e\,\AA$^{-2}$ and $s$ the scan step. Doses
of $10^{4}$--$10^{10}$\,e\,\AA$^{-2}$ are used: the coherent, high-fluence limit furnishes
the near-noiseless reference reconstruction, and the thermally-averaged, counting-limited
data the experimentally realistic test.

\subsection{Multislice electron ptychography}
\label{sec:recon}

\paragraph{Ptychographic inversion.}
The complex specimen transmission is recovered by iterative ptychography
\citep{RodenburgFaulkner2004,Wakonig2020}, which exploits the redundancy of overlapping
coherent diffraction patterns and has reached deep sub-\aa{}ngstr\"om lateral resolution for
electrons \citep{Jiang2018}. To resolve structure \emph{along} the beam the single
multiplicative projection is replaced by a multislice object \citep{MaidenHumphryRodenburg2012}:
the specimen is modelled as a stack of $N_L$ transmission slices with propagation between
them, exactly as in Eq.~\eqref{eq:multislice}, and the inversion jointly recovers all
slices. This is what converts a single-orientation ptychographic measurement into a
depth-resolved, optically sectioned reconstruction without a tilt series
\citep{Gao2017,Chen2021}. We reconstruct $N_L=70$ (coherent) or $105$ (thermally-averaged)
slices across the $\approx 70$\,\AA{} thickness, a depth sampling of $\Delta z_\mathrm{rec}
\approx 1.0$ and $0.67$\,\AA{} respectively, with the hollow-detector inner angle set to
zero so that the full recorded cone constrains the object.

\paragraph{Maximum-likelihood solver.}
Reconstruction uses the least-squares maximum-likelihood (LSQ-ML) engine of
\citet{Odstrcil2018} within PtychoShelves \citep{Wakonig2020}, which descends a
likelihood-based data-fidelity objective \citep{ThibaultGuizar2012} by projected
least-squares updates of the object slices and probe. A two-stage schedule is used: a
down-sampled pre-solve on $178$-pixel diffraction arrays fixes the low-frequency object and
probe before a full-resolution $356$-pixel refinement, which accelerates convergence and
reduces the risk of stagnation in a poor local minimum. Inter-layer regularisation is
disabled, so the depth distribution is set by the data rather than by a smoothness prior;
the cost is a noisier deep reconstruction, which we characterise directly rather than
suppress (Sec.~\ref{sec:atomfind}).

\paragraph{Partial coherence and the probe.}
For the coherent reference a single, fixed probe is used. For the thermally-averaged,
counting-limited data the illumination is described by a mixed state --- an incoherent sum
of mutually orthogonal probe modes $\{P_m\}$ \citep{ThibaultMenzel2013},
\begin{equation}
  I(\mathbf{k}) \;=\; \sum_{m} \big|\,\mathcal{F}\{P_m\, O\}(\mathbf{k})\,\big|^2 ,
  \label{eq:mixed}
\end{equation}
released after a single-mode burn-in. The additional incoherent degrees of freedom absorb
the thermal diffuse background and residual partial coherence that a single-mode model would
otherwise be forced to explain with the coherent object, corrupting it
\citep{ThibaultMenzel2013}; slow probe variation across the scan, where present, is further
accommodated by orthogonal probe relaxation \citep{Odstrcil2016}.

\paragraph{Depth resolution and the missing cone.}
Because a single-orientation measurement samples reciprocal space only within the aperture
cone, the axial resolution is bounded by the numerical aperture as $\approx\lambda/\alpha^2$
($\approx 2$\,\AA{} for this probe) against a lateral resolution $\approx\lambda/2\alpha$
($\approx 0.1$\,\AA); this is the missing-cone, limited-angle regime of depth recovery,
here reached by multislice electron ptychography rather than by tomographic tilting
\citep{Chen2021}. The reconstructed response of a single atom inherits exactly this
anisotropy and is measured directly to serve as the three-dimensional point-spread function
of the localisation model in Sec.~\ref{sec:atomfind}.
% ---------------------------------------------------------------------------
